% --- Archon physics formalization source begin ---
% archon:physics
% archon:covers PhyXMiniProblems/problem_phyx_mini_0354.lean
% archon:source-report reports/phyx_mini/problem_phyx_mini_0354.source.json

\chapter{Physics problem phyx\_mini\_0354}
\label{ch:PhyXMiniProblems_problem_phyx_mini_0354}

\paragraph{Problem source.}
\#\# Physical scenario

You are conducting experiments to study prototype heat engines. In one test, 4.00 mol of argon gas are taken around the cycle shown in figure. The pressure is low enough for the gas to be treated as ideal. You measure the gas temperature in states \$a\$, \$b\$, \$c\$, and \$d\$ and find \$T\_a = 250.0\$ K, \$T\_b = 300.0\$ K, \$T\_c = 380.0\$ K, and \$T\_d = 316.7\$ K.

\#\# Question

Calculate the efficiency \$e\$ of the cycle.

\#\# Answer choices

- A: A: 10\textbackslash{}\%

- B: B: 6.25\textbackslash{}\%

- C: C: 5.9\textbackslash{}\%

- D: D: 12.3\textbackslash{}\%

\#\# Figure caption (auxiliary; use the image as primary evidence)

The image depicts a PV (pressure-volume) diagram with a cycle process:

- **Axes:** The vertical axis is labeled \textbackslash{}( p \textbackslash{}) and the horizontal axis is labeled \textbackslash{}( V \textbackslash{}), representing pressure and volume, respectively. The origin is labeled as \textbackslash{}( O \textbackslash{}).

- **Cycle:** A rectangular cycle is shown, consisting of four points labeled \textbackslash{}( a \textbackslash{}), \textbackslash{}( b \textbackslash{}), \textbackslash{}( c \textbackslash{}), and \textbackslash{}( d \textbackslash{}).

- **Lines and Arrows:** The cycle is made up of four lines connecting these points:
  - Line from \textbackslash{}( a \textbackslash{}) to \textbackslash{}( b \textbackslash{})
  - Line from \textbackslash{}( b \textbackslash{}) to \textbackslash{}( c \textbackslash{})
  - Line from \textbackslash{}( c \textbackslash{}) to \textbackslash{}( d \textbackslash{})
  - Line from \textbackslash{}( d \textbackslash{}) to \textbackslash{}( a \textbackslash{})

  The lines have arrows indicating the direction of the cycle:
  - \textbackslash{}( a \textbackslash{}rightarrow b \textbackslash{}rightarrow c \textbackslash{}rightarrow d \textbackslash{}rightarrow a \textbackslash{})

- **Color:** The lines and arrows are colored blue.

- **Points:** Each point \textbackslash{}( a \textbackslash{}), \textbackslash{}( b \textbackslash{}), \textbackslash{}( c \textbackslash{}), \textbackslash{}( d \textbackslash{}) is marked with a black dot.

This diagram typically represents a thermodynamic cycle, where processes between the points involve changes in pressure and volume.

\#\# Dataset metadata

Laws of Thermodynamics; Physical Model Grounding Reasoning, Multi-Formula Reasoning

\paragraph{Recorded answer/context.}
B: B: 6.25\textbackslash{}\%

\paragraph{Figure/image path.}
/root/proposal\_for\_physic/science-mango/phyx\_mini\_run/phyx\_data/test\_image/354.png

\paragraph{Formalization target.}
create a compiling Lean file with sorry bodies at `PhyXMiniProblems/problem\_phyx\_mini\_0354.lean`. The Lean declarations must preserve the physical quantities, dimensions or dimensional roles, figure labels, governing-law hypotheses, and final relation expressed by this problem.
Use Mathlib/Physlib names found through LeanExplore where available. If a physics API is missing, introduce faithful local abstractions rather than scalar placeholder aliases.

\begin{theorem}[Physics formalization target]
\label{thm:physics:phyx_mini_0354:target}
\lean{PhyXMiniProblems.ProblemPhyXMini0354.problem_phyx_mini_0354}
\uses{def:physics:phyx-mini-0354:phyxminiproblems-problemphyxmini0354-matchessuppliedrectangularcycle, def:physics:phyx-mini-0354:phyxminiproblems-problemphyxmini0354-argonheatenginesetup, def:physics:phyx-mini-0354:phyxminiproblems-problemphyxmini0354-matchesargonsampleandtemperaturereadouts, def:physics:phyx-mini-0354:phyxminiproblems-problemphyxmini0354-hasphysicalcycleparameters, def:physics:phyx-mini-0354:phyxminiproblems-problemphyxmini0354-satisfiesmonatomicidealgascyclelaws, def:physics:phyx-mini-0354:phyxminiproblems-problemphyxmini0354-thermalefficiency, def:physics:phyx-mini-0354:phyxminiproblems-problemphyxmini0354-thermalefficiencypercent, def:physics:phyx-mini-0354:phyxminiproblems-problemphyxmini0354-matchesestimate, def:physics:phyx-mini-0354:phyxminiproblems-problemphyxmini0354-iscompatibledisplayedefficiencyanswer}
The assigned autoformalize agent should translate this physics problem into Lean declarations in the covered file, with theorem and lemma proof bodies written as `by sorry`.
\end{theorem}
\begin{proof}
This is an autoformalization task, not a proof task. Produce statements that can later be proved by the physics prover without weakening the physical model.
\end{proof}

% --- Archon declaration topology begin ---
\section{Declaration topology}
\begin{definition}[Volume Quantity]
  \label{def:physics:phyx-mini-0354:phyxminiproblems-problemphyxmini0354-volumequantity}
  \lean{PhyXMiniProblems.ProblemPhyXMini0354.VolumeQuantity}
  A nonnegative physical volume with dimension length cubed
\end{definition}
\begin{proof}
  This is the declared model definition of Volume Quantity.
\end{proof}
\begin{definition}[volume In Cubic Meters]
  \label{def:physics:phyx-mini-0354:phyxminiproblems-problemphyxmini0354-volumeincubicmeters}
  \lean{PhyXMiniProblems.ProblemPhyXMini0354.volumeInCubicMeters}
  \uses{def:physics:phyx-mini-0354:phyxminiproblems-problemphyxmini0354-volumequantity}
  Cubic-metre readout of a physical volume
\end{definition}
\begin{proof}
  This is the declared model definition of volume In Cubic Meters.
\end{proof}
\begin{definition}[pressure In Pascals]
  \label{def:physics:phyx-mini-0354:phyxminiproblems-problemphyxmini0354-pressureinpascals}
  \lean{PhyXMiniProblems.ProblemPhyXMini0354.pressureInPascals}
  Pascal readout of a physical pressure
\end{definition}
\begin{proof}
  This is the declared model definition of pressure In Pascals.
\end{proof}
\begin{definition}[energy In Joules]
  \label{def:physics:phyx-mini-0354:phyxminiproblems-problemphyxmini0354-energyinjoules}
  \lean{PhyXMiniProblems.ProblemPhyXMini0354.energyInJoules}
  Joule readout of a physical energy, heat transfer, or work
\end{definition}
\begin{proof}
  This is the declared model definition of energy In Joules.
\end{proof}
\begin{definition}[temperature In Kelvin]
  \label{def:physics:phyx-mini-0354:phyxminiproblems-problemphyxmini0354-temperatureinkelvin}
  \lean{PhyXMiniProblems.ProblemPhyXMini0354.temperatureInKelvin}
  Kelvin readout of a Physlib absolute temperature stored in an explicitly chosen zero-preserving temperature unit
\end{definition}
\begin{proof}
  This is the declared model definition of temperature In Kelvin.
\end{proof}
\begin{definition}[Cycle State]
  \label{def:physics:phyx-mini-0354:phyxminiproblems-problemphyxmini0354-cyclestate}
  \lean{PhyXMiniProblems.ProblemPhyXMini0354.CycleState}
  The four black vertices labelled in the supplied pressure--volume plot
\end{definition}
\begin{proof}
  This is the declared model definition of Cycle State.
\end{proof}
\begin{definition}[Cycle Leg]
  \label{def:physics:phyx-mini-0354:phyxminiproblems-problemphyxmini0354-cycleleg}
  \lean{PhyXMiniProblems.ProblemPhyXMini0354.CycleLeg}
  The four directed blue sides of the plotted cycle
\end{definition}
\begin{proof}
  This is the declared model definition of Cycle Leg.
\end{proof}
\begin{definition}[Leg Constraint]
  \label{def:physics:phyx-mini-0354:phyxminiproblems-problemphyxmini0354-legconstraint}
  \lean{PhyXMiniProblems.ProblemPhyXMini0354.LegConstraint}
  The thermodynamic constraint represented by a straight plot side
\end{definition}
\begin{proof}
  This is the declared model definition of Leg Constraint.
\end{proof}
\begin{definition}[Diagram Label]
  \label{def:physics:phyx-mini-0354:phyxminiproblems-problemphyxmini0354-diagramlabel}
  \lean{PhyXMiniProblems.ProblemPhyXMini0354.DiagramLabel}
  Text labels explicitly present in the primary image
\end{definition}
\begin{proof}
  This is the declared model definition of Diagram Label.
\end{proof}
\begin{definition}[leg Source]
  \label{def:physics:phyx-mini-0354:phyxminiproblems-problemphyxmini0354-legsource}
  \lean{PhyXMiniProblems.ProblemPhyXMini0354.legSource}
  \uses{def:physics:phyx-mini-0354:phyxminiproblems-problemphyxmini0354-cyclestate, def:physics:phyx-mini-0354:phyxminiproblems-problemphyxmini0354-cycleleg}
  Source state of a directed cycle leg
\end{definition}
\begin{proof}
  This is the declared model definition of leg Source.
\end{proof}
\begin{definition}[leg Target]
  \label{def:physics:phyx-mini-0354:phyxminiproblems-problemphyxmini0354-legtarget}
  \lean{PhyXMiniProblems.ProblemPhyXMini0354.legTarget}
  \uses{def:physics:phyx-mini-0354:phyxminiproblems-problemphyxmini0354-cyclestate, def:physics:phyx-mini-0354:phyxminiproblems-problemphyxmini0354-cycleleg}
  Target state of a directed cycle leg
\end{definition}
\begin{proof}
  This is the declared model definition of leg Target.
\end{proof}
\begin{definition}[displayed Leg Constraint]
  \label{def:physics:phyx-mini-0354:phyxminiproblems-problemphyxmini0354-displayedlegconstraint}
  \lean{PhyXMiniProblems.ProblemPhyXMini0354.displayedLegConstraint}
  \uses{def:physics:phyx-mini-0354:phyxminiproblems-problemphyxmini0354-cycleleg, def:physics:phyx-mini-0354:phyxminiproblems-problemphyxmini0354-legconstraint}
  Vertical sides are isochoric and horizontal sides are isobaric
\end{definition}
\begin{proof}
  This is the declared model definition of displayed Leg Constraint.
\end{proof}
\begin{definition}[Pressure Volume Cycle Figure]
  \label{def:physics:phyx-mini-0354:phyxminiproblems-problemphyxmini0354-pressurevolumecyclefigure}
  \lean{PhyXMiniProblems.ProblemPhyXMini0354.PressureVolumeCycleFigure}
  \uses{def:physics:phyx-mini-0354:phyxminiproblems-problemphyxmini0354-volumequantity, def:physics:phyx-mini-0354:phyxminiproblems-problemphyxmini0354-cyclestate, def:physics:phyx-mini-0354:phyxminiproblems-problemphyxmini0354-cycleleg, def:physics:phyx-mini-0354:phyxminiproblems-problemphyxmini0354-legconstraint, def:physics:phyx-mini-0354:phyxminiproblems-problemphyxmini0354-diagramlabel}
  Physical coordinates and qualitative marks in the supplied diagram
\end{definition}
\begin{proof}
  This is the declared model definition of Pressure Volume Cycle Figure.
\end{proof}
\begin{definition}[Matches Supplied Rectangular Cycle]
  \label{def:physics:phyx-mini-0354:phyxminiproblems-problemphyxmini0354-matchessuppliedrectangularcycle}
  \lean{PhyXMiniProblems.ProblemPhyXMini0354.MatchesSuppliedRectangularCycle}
  \uses{def:physics:phyx-mini-0354:phyxminiproblems-problemphyxmini0354-volumeincubicmeters, def:physics:phyx-mini-0354:phyxminiproblems-problemphyxmini0354-pressureinpascals, def:physics:phyx-mini-0354:phyxminiproblems-problemphyxmini0354-displayedlegconstraint, def:physics:phyx-mini-0354:phyxminiproblems-problemphyxmini0354-pressurevolumecyclefigure}
  Primary-image evidence for the four labelled vertices, directed arrows, and equal-coordinate and ordering relations of the rectangular cycle
\end{definition}
\begin{proof}
  This is the declared model definition of Matches Supplied Rectangular Cycle.
\end{proof}
\begin{definition}[Gas Species]
  \label{def:physics:phyx-mini-0354:phyxminiproblems-problemphyxmini0354-gasspecies}
  \lean{PhyXMiniProblems.ProblemPhyXMini0354.GasSpecies}
  Chemical species named in the experiment
\end{definition}
\begin{proof}
  This is the declared model definition of Gas Species.
\end{proof}
\begin{definition}[Gas Model]
  \label{def:physics:phyx-mini-0354:phyxminiproblems-problemphyxmini0354-gasmodel}
  \lean{PhyXMiniProblems.ProblemPhyXMini0354.GasModel}
  Equation-of-state and caloric model assigned to the gas
\end{definition}
\begin{proof}
  This is the declared model definition of Gas Model.
\end{proof}
\begin{definition}[Argon Heat Engine Setup]
  \label{def:physics:phyx-mini-0354:phyxminiproblems-problemphyxmini0354-argonheatenginesetup}
  \lean{PhyXMiniProblems.ProblemPhyXMini0354.ArgonHeatEngineSetup}
  \uses{def:physics:phyx-mini-0354:phyxminiproblems-problemphyxmini0354-cyclestate, def:physics:phyx-mini-0354:phyxminiproblems-problemphyxmini0354-cycleleg, def:physics:phyx-mini-0354:phyxminiproblems-problemphyxmini0354-pressurevolumecyclefigure, def:physics:phyx-mini-0354:phyxminiproblems-problemphyxmini0354-gasspecies, def:physics:phyx-mini-0354:phyxminiproblems-problemphyxmini0354-gasmodel}
  Physical data for the engine. The amount of substance remains abstract and only its mole readout is real-valued. Heat is positive into the gas, while work is positive when done by the gas
\end{definition}
\begin{proof}
  This is the declared model definition of Argon Heat Engine Setup.
\end{proof}
\begin{definition}[mole Amount Of Gas]
  \label{def:physics:phyx-mini-0354:phyxminiproblems-problemphyxmini0354-moleamountofgas}
  \lean{PhyXMiniProblems.ProblemPhyXMini0354.moleAmountOfGas}
  \uses{def:physics:phyx-mini-0354:phyxminiproblems-problemphyxmini0354-argonheatenginesetup}
  Mole readout of the physical argon sample
\end{definition}
\begin{proof}
  This is the declared model definition of mole Amount Of Gas.
\end{proof}
\begin{definition}[gas Temperature In Kelvin]
  \label{def:physics:phyx-mini-0354:phyxminiproblems-problemphyxmini0354-gastemperatureinkelvin}
  \lean{PhyXMiniProblems.ProblemPhyXMini0354.gasTemperatureInKelvin}
  \uses{def:physics:phyx-mini-0354:phyxminiproblems-problemphyxmini0354-temperatureinkelvin, def:physics:phyx-mini-0354:phyxminiproblems-problemphyxmini0354-cyclestate, def:physics:phyx-mini-0354:phyxminiproblems-problemphyxmini0354-argonheatenginesetup}
  Kelvin readout at a labelled cycle state
\end{definition}
\begin{proof}
  This is the declared model definition of gas Temperature In Kelvin.
\end{proof}
\begin{definition}[internal Energy In Joules]
  \label{def:physics:phyx-mini-0354:phyxminiproblems-problemphyxmini0354-internalenergyinjoules}
  \lean{PhyXMiniProblems.ProblemPhyXMini0354.internalEnergyInJoules}
  \uses{def:physics:phyx-mini-0354:phyxminiproblems-problemphyxmini0354-energyinjoules, def:physics:phyx-mini-0354:phyxminiproblems-problemphyxmini0354-cyclestate, def:physics:phyx-mini-0354:phyxminiproblems-problemphyxmini0354-argonheatenginesetup}
  Joule readout of internal energy at a labelled state
\end{definition}
\begin{proof}
  This is the declared model definition of internal Energy In Joules.
\end{proof}
\begin{definition}[work Done By Gas In Joules]
  \label{def:physics:phyx-mini-0354:phyxminiproblems-problemphyxmini0354-workdonebygasinjoules}
  \lean{PhyXMiniProblems.ProblemPhyXMini0354.workDoneByGasInJoules}
  \uses{def:physics:phyx-mini-0354:phyxminiproblems-problemphyxmini0354-energyinjoules, def:physics:phyx-mini-0354:phyxminiproblems-problemphyxmini0354-cycleleg, def:physics:phyx-mini-0354:phyxminiproblems-problemphyxmini0354-argonheatenginesetup}
  Signed joule readout of work done by the gas along one directed leg
\end{definition}
\begin{proof}
  This is the declared model definition of work Done By Gas In Joules.
\end{proof}
\begin{definition}[heat Transferred Into Gas In Joules]
  \label{def:physics:phyx-mini-0354:phyxminiproblems-problemphyxmini0354-heattransferredintogasinjoules}
  \lean{PhyXMiniProblems.ProblemPhyXMini0354.heatTransferredIntoGasInJoules}
  \uses{def:physics:phyx-mini-0354:phyxminiproblems-problemphyxmini0354-energyinjoules, def:physics:phyx-mini-0354:phyxminiproblems-problemphyxmini0354-cycleleg, def:physics:phyx-mini-0354:phyxminiproblems-problemphyxmini0354-argonheatenginesetup}
  Signed joule readout of heat transferred into the gas along one leg
\end{definition}
\begin{proof}
  This is the declared model definition of heat Transferred Into Gas In Joules.
\end{proof}
\begin{definition}[temperature Readout Tolerance In Kelvin]
  \label{def:physics:phyx-mini-0354:phyxminiproblems-problemphyxmini0354-temperaturereadouttoleranceinkelvin}
  \lean{PhyXMiniProblems.ProblemPhyXMini0354.temperatureReadoutToleranceInKelvin}
  Half the final reported decimal place, in kelvin
\end{definition}
\begin{proof}
  This is the declared model definition of temperature Readout Tolerance In Kelvin.
\end{proof}
\begin{definition}[Matches Temperature Readout]
  \label{def:physics:phyx-mini-0354:phyxminiproblems-problemphyxmini0354-matchestemperaturereadout}
  \lean{PhyXMiniProblems.ProblemPhyXMini0354.MatchesTemperatureReadout}
  \uses{def:physics:phyx-mini-0354:phyxminiproblems-problemphyxmini0354-temperaturereadouttoleranceinkelvin}
  Compatibility with a temperature reported to one decimal place
\end{definition}
\begin{proof}
  This is the declared model definition of Matches Temperature Readout.
\end{proof}
\begin{definition}[Matches Argon Sample And Temperature Readouts]
  \label{def:physics:phyx-mini-0354:phyxminiproblems-problemphyxmini0354-matchesargonsampleandtemperaturereadouts}
  \lean{PhyXMiniProblems.ProblemPhyXMini0354.MatchesArgonSampleAndTemperatureReadouts}
  \uses{def:physics:phyx-mini-0354:phyxminiproblems-problemphyxmini0354-argonheatenginesetup, def:physics:phyx-mini-0354:phyxminiproblems-problemphyxmini0354-moleamountofgas, def:physics:phyx-mini-0354:phyxminiproblems-problemphyxmini0354-gastemperatureinkelvin, def:physics:phyx-mini-0354:phyxminiproblems-problemphyxmini0354-matchestemperaturereadout}
  Problem-prose data: an ideal monatomic four-mole argon sample and the four measured temperature readouts. No heat, work, or efficiency is fixed here
\end{definition}
\begin{proof}
  This is the declared model definition of Matches Argon Sample And Temperature Readouts.
\end{proof}
\begin{definition}[Has Physical Cycle Parameters]
  \label{def:physics:phyx-mini-0354:phyxminiproblems-problemphyxmini0354-hasphysicalcycleparameters}
  \lean{PhyXMiniProblems.ProblemPhyXMini0354.HasPhysicalCycleParameters}
  \uses{def:physics:phyx-mini-0354:phyxminiproblems-problemphyxmini0354-volumeincubicmeters, def:physics:phyx-mini-0354:phyxminiproblems-problemphyxmini0354-pressureinpascals, def:physics:phyx-mini-0354:phyxminiproblems-problemphyxmini0354-argonheatenginesetup, def:physics:phyx-mini-0354:phyxminiproblems-problemphyxmini0354-moleamountofgas, def:physics:phyx-mini-0354:phyxminiproblems-problemphyxmini0354-gastemperatureinkelvin}
  Positivity and nondegeneracy conditions for a physical engine cycle
\end{definition}
\begin{proof}
  This is the declared model definition of Has Physical Cycle Parameters.
\end{proof}
\begin{definition}[Satisfies Monatomic Ideal Gas Cycle Laws]
  \label{def:physics:phyx-mini-0354:phyxminiproblems-problemphyxmini0354-satisfiesmonatomicidealgascyclelaws}
  \lean{PhyXMiniProblems.ProblemPhyXMini0354.SatisfiesMonatomicIdealGasCycleLaws}
  \uses{def:physics:phyx-mini-0354:phyxminiproblems-problemphyxmini0354-volumeincubicmeters, def:physics:phyx-mini-0354:phyxminiproblems-problemphyxmini0354-pressureinpascals, def:physics:phyx-mini-0354:phyxminiproblems-problemphyxmini0354-legsource, def:physics:phyx-mini-0354:phyxminiproblems-problemphyxmini0354-legtarget, def:physics:phyx-mini-0354:phyxminiproblems-problemphyxmini0354-argonheatenginesetup, def:physics:phyx-mini-0354:phyxminiproblems-problemphyxmini0354-moleamountofgas, def:physics:phyx-mini-0354:phyxminiproblems-problemphyxmini0354-gastemperatureinkelvin, def:physics:phyx-mini-0354:phyxminiproblems-problemphyxmini0354-internalenergyinjoules, def:physics:phyx-mini-0354:phyxminiproblems-problemphyxmini0354-workdonebygasinjoules, def:physics:phyx-mini-0354:phyxminiproblems-problemphyxmini0354-heattransferredintogasinjoules}
  Thermodynamic laws used for this low-pressure monatomic ideal-gas cycle: p V = n R T, U = (3/2) n R T, quasistatic boundary work, and the first law Q = ΔU + W with the sign conventions declared above
\end{definition}
\begin{proof}
  This is the declared model definition of Satisfies Monatomic Ideal Gas Cycle Laws.
\end{proof}
\begin{definition}[net Cycle Work In Joules]
  \label{def:physics:phyx-mini-0354:phyxminiproblems-problemphyxmini0354-netcycleworkinjoules}
  \lean{PhyXMiniProblems.ProblemPhyXMini0354.netCycleWorkInJoules}
  \uses{def:physics:phyx-mini-0354:phyxminiproblems-problemphyxmini0354-cycleleg, def:physics:phyx-mini-0354:phyxminiproblems-problemphyxmini0354-argonheatenginesetup, def:physics:phyx-mini-0354:phyxminiproblems-problemphyxmini0354-workdonebygasinjoules}
  Net work done by the gas over the complete directed cycle, in joules
\end{definition}
\begin{proof}
  This is the declared model definition of net Cycle Work In Joules.
\end{proof}
\begin{definition}[total Heat Input In Joules]
  \label{def:physics:phyx-mini-0354:phyxminiproblems-problemphyxmini0354-totalheatinputinjoules}
  \lean{PhyXMiniProblems.ProblemPhyXMini0354.totalHeatInputInJoules}
  \uses{def:physics:phyx-mini-0354:phyxminiproblems-problemphyxmini0354-cycleleg, def:physics:phyx-mini-0354:phyxminiproblems-problemphyxmini0354-argonheatenginesetup, def:physics:phyx-mini-0354:phyxminiproblems-problemphyxmini0354-heattransferredintogasinjoules}
  Sum of the positive parts of the signed heat transfers, in joules
\end{definition}
\begin{proof}
  This is the declared model definition of total Heat Input In Joules.
\end{proof}
\begin{definition}[thermal Efficiency]
  \label{def:physics:phyx-mini-0354:phyxminiproblems-problemphyxmini0354-thermalefficiency}
  \lean{PhyXMiniProblems.ProblemPhyXMini0354.thermalEfficiency}
  \uses{def:physics:phyx-mini-0354:phyxminiproblems-problemphyxmini0354-argonheatenginesetup, def:physics:phyx-mini-0354:phyxminiproblems-problemphyxmini0354-netcycleworkinjoules, def:physics:phyx-mini-0354:phyxminiproblems-problemphyxmini0354-totalheatinputinjoules}
  Dimensionless thermal efficiency W\_net / Q\_in
\end{definition}
\begin{proof}
  This is the declared model definition of thermal Efficiency.
\end{proof}
\begin{definition}[thermal Efficiency Percent]
  \label{def:physics:phyx-mini-0354:phyxminiproblems-problemphyxmini0354-thermalefficiencypercent}
  \lean{PhyXMiniProblems.ProblemPhyXMini0354.thermalEfficiencyPercent}
  \uses{def:physics:phyx-mini-0354:phyxminiproblems-problemphyxmini0354-argonheatenginesetup, def:physics:phyx-mini-0354:phyxminiproblems-problemphyxmini0354-thermalefficiency}
  Thermal efficiency expressed as a percentage
\end{definition}
\begin{proof}
  This is the declared model definition of thermal Efficiency Percent.
\end{proof}
\begin{definition}[Answer Choice]
  \label{def:physics:phyx-mini-0354:phyxminiproblems-problemphyxmini0354-answerchoice}
  \lean{PhyXMiniProblems.ProblemPhyXMini0354.AnswerChoice}
  Labels of the four printed multiple-choice answers
\end{definition}
\begin{proof}
  This is the declared model definition of Answer Choice.
\end{proof}
\begin{definition}[displayed Efficiency Percent]
  \label{def:physics:phyx-mini-0354:phyxminiproblems-problemphyxmini0354-displayedefficiencypercent}
  \lean{PhyXMiniProblems.ProblemPhyXMini0354.displayedEfficiencyPercent}
  \uses{def:physics:phyx-mini-0354:phyxminiproblems-problemphyxmini0354-answerchoice}
  Percentage printed beside each answer label
\end{definition}
\begin{proof}
  This is the declared model definition of displayed Efficiency Percent.
\end{proof}
\begin{definition}[recorded Answer Choice]
  \label{def:physics:phyx-mini-0354:phyxminiproblems-problemphyxmini0354-recordedanswerchoice}
  \lean{PhyXMiniProblems.ProblemPhyXMini0354.recordedAnswerChoice}
  \uses{def:physics:phyx-mini-0354:phyxminiproblems-problemphyxmini0354-answerchoice}
  The answer label recorded by the dataset, retained as source metadata
\end{definition}
\begin{proof}
  This is the declared model definition of recorded Answer Choice.
\end{proof}
\begin{definition}[Matches Estimate]
  \label{def:physics:phyx-mini-0354:phyxminiproblems-problemphyxmini0354-matchesestimate}
  \lean{PhyXMiniProblems.ProblemPhyXMini0354.MatchesEstimate}
  An actual scalar is within tolerance of a stated estimate
\end{definition}
\begin{proof}
  This is the declared model definition of Matches Estimate.
\end{proof}
\begin{definition}[Is Compatible Displayed Efficiency Answer]
  \label{def:physics:phyx-mini-0354:phyxminiproblems-problemphyxmini0354-iscompatibledisplayedefficiencyanswer}
  \lean{PhyXMiniProblems.ProblemPhyXMini0354.IsCompatibleDisplayedEfficiencyAnswer}
  \uses{def:physics:phyx-mini-0354:phyxminiproblems-problemphyxmini0354-answerchoice, def:physics:phyx-mini-0354:phyxminiproblems-problemphyxmini0354-displayedefficiencypercent, def:physics:phyx-mini-0354:phyxminiproblems-problemphyxmini0354-matchesestimate}
  A computed percentage is compatible with a displayed choice to one tenth of a percentage point
\end{definition}
\begin{proof}
  This is the declared model definition of Is Compatible Displayed Efficiency Answer.
\end{proof}
% --- Archon declaration topology end ---
% --- Archon physics formalization source end ---
